\documentclass{article}
\usepackage{../math_template}

\author{Jonathan Kasongo}
\title{Mathematical Excalibur --- Volume 7, No 3}

\begin{document}
\maketitle

\begin{problem}{Example 1.}
Numbers $1, 2, 3, \ldots, 1974$ are written on
a board. You are allowed to replace any
two of these numbers by one number,
which is either the sum or the difference
of these numbers. Show that after 1973
times performing this operation, the only
number left on the board cannot be 0.
\\

\textit{(1974 Kiev Math Olympiad)}
\end{problem}

\begin{solution}{Solution to Example 1.}
Initially there are 987 odd numbers and 987 even numbers.

\begin{itemize}
\item If we choose to replace two numbers of the same parity, their sum and
difference is even, so we either lose two odd numbers or we don't lose any
odd numbers.
\item If we choose to replace two numbers of different parity, their
sum and difference is odd, so the number of odd numbers in the list
doesn't change.
\end{itemize}

In each of these cases, the parity of the number of odd numbers on the board
doesn't change and so it is an invariant. Thus the last number left on the
board must be odd, and so it cannot be 0. $\blacksquare$
\end{solution}
\vspace{0.2cm}

\begin{problem}{Example 2.}
In an $8 \times 8$ board, there are
32 white pieces and 32 black pieces, one
piece in each square. If a player can
change all the white pieces to black and
all the black pieces to white in any row
or column in a single move, then is it
possible that after finitely many moves,
there will be exactly one black piece left
on the board?
\end{problem}

\begin{solution}{Solution to Example 2.}
Let's say that in one row or column, there are $w$ white pieces and $b$
black pieces. Indeed $w + b = 8 \equiv 0 \ \text{(mod } 2)$. Thus
$w \equiv b \ \text{(mod } 2)$. The total number of black pieces on the
board will increase by $w-b$ each move, and similarly the total number of
white pieces on the board will increase by $b-w$ each move. Since both
$b-w$ and $w-b$ have to be even, the number of white pieces on the board at
any given time is even, and the same for the number of black pieces on the
board at any given time. Thus there can never only be 1 black piece left
on the board. $\blacksquare$
\end{solution}
\vspace{0.2cm}

\begin{problem}{Example 3.}
Four $x$’s and five $o$’s are
written around the circle in an arbitrary
order. If two consecutive symbols are the
same, then insert a new $x$ between them.
Otherwise insert a new $o$ between them.
Remove the old $x$’s and $o$’s. Keep on
repeating this operation. Is it possible to
get nine $o$’s?
\end{problem}

\begin{solution}{Solution to Example 3.}
\begin{itemize}
\item If we have two consecutive $x$'s, then we will replace them with a
singular $x$, and the number of $o$'s doesn't change.
\item If we have two consecutive $o$'s then we will replace them with a
singular $x$, and the number of $o$'s will decrease by 2.
\item If we have an $x$ next to an $o$ then we will replace them with a
singular $o$, and the number of $o$'s doesn't change.
\end{itemize}

In each case the number of $o$'s decreases and so it is impossible to get
nine $o$'s. $\blacksquare$
\end{solution}
\vspace{0.2cm}

\begin{problem}{Example 4.}
There are three piles of
stones numbering 19, 8 and 9,
respectively. You are allowed to choose
two piles and transfer one stone from
each of these two piles to the third piles.
After several of these operations, is it
possible that each of the three piles has
12 stones?
\end{problem}

\begin{solution}{Solution to Example 4.}
Modulo 3 the piles of stones are 1, 2 and 0 respectively. If a pile receives
two stones, or if it has one stone removed this corresponds to adding
$2$ to the pile modulo 3. Since after each operation each of the piles
modulo 3 change by the same amount, we can never have them all be equal
to 0 mod 3 at some point, thus it is never possible to have each of the
three piles with 12 stones. $\blacksquare$
\\

\textit{Remark: The motivation for mod 3, comes from the fact that $2
\equiv -1 \ (\text{mod } 3)$.}
\end{solution}
\end{document}
